\documentclass[10pt,a4paper]{amsart}
\usepackage[margin=19mm]{geometry}
\usepackage{amsmath,amssymb,mathtools}
\usepackage[scr=rsfs,cal=euler]{mathalfa}
\usepackage{xfrac}
\usepackage[unicode,hidelinks,bookmarks=false]{hyperref}
\newcommand{\ie}{i.e.\ }
\newcommand{\cf}{cf.\ }
\newcommand{\resp}{respectively\ }
\newcommand{\Subsection}{\subsection}
\providecommand{\nobreakdash}{\nobreak\hbox{-}}
\setlength{\emergencystretch}{2em}
\multlinegap0pt
\usepackage{bm}
\usepackage{bbm}
\usepackage{dsfont}
\usepackage{xspace}
\usepackage{cancel}
\usepackage{mathtools}
\usepackage{amsthm}
\makeatletter
\@ifundefined{theorem}{\newtheorem{theorem}{Theorem}}{}
\@ifundefined{lemma}{\newtheorem{lemma}[theorem]{Lemma}}{}
\makeatother
\DeclareMathOperator{\supp}{supp}
\title[Fractional Quantum Hall States: Infinite Matrix Product Repr]{Fractional Quantum Hall States: Infinite Matrix Product Representation and its Implications}
\author{Severin Schraven, Simone Warzel}
\date{Source: arXiv:2601.12165v1 (CC BY 4.0)}
\begin{document}
\maketitle

\noindent\emph{Source and licence.} Fractional Quantum Hall States: Infinite Matrix Product Representation and its Implications. Version arXiv:2601.12165v1 (CC BY 4.0), distributed under the Creative Commons Attribution 4.0 International licence, as recorded in the arXiv licence metadata for this exact version and shown on the abstract page https://arxiv.org/abs/2601.12165v1. Full rights notice: licenses/arxiv-2601.12165-CC-BY-4.0.txt.\par\smallskip

\noindent\textit{Editorial scope.} Counted unit: the Lemma whose source label is \texttt{lm:partitions PsiAB} in the article (source0.tex lines 1778--1789, source label \texttt{lm:partitions PsiAB}), delivered together with the article's own complete original proof of it (source0.tex lines 1790--1855, one \texttt{proof} environment of 66 lines). No mathematics is written, repaired or re-derived: statement and proof are the article's own text line for line. Required context retained uncounted: none. Every definition, display and label of those lines is retained unchanged. Omitted from this package and not redistributed: 1--1777, 1856--2309. Nothing inside a retained excerpt is omitted or altered: every retained body is delivered exactly as the article's lines are written, so no omission receipt of any kind accompanies this package. Citations: the retained excerpts carry no citation command at all, so no bibliography entry of the article is transcribed and no reference is redistributed. Reference adaptation: none was needed and none was made; every reference inside the delivered unit resolves inside it, so no unresolved reference or citation is produced. Printed slips kept literal: no printed word, symbol, hypothesis or number of the counted statement or of its proof was altered. Layout and class: the article's own class is \texttt{scrartcl} with packages including amssymb, amsthm, mathtools, bbm, bm, natbib, hyphenat, babel, enumerate, amsmath, dsfont, hyperref, color, inputenc; the wrapper is the lane's pinned \texttt{amsart} editorial wrapper at 10pt on A4, which restates the theorem-like environments the retained text needs and adds the article's own macro definitions that the retained text needs (\texttt{\textbackslash supp}), with extra wrapper packages (bm, bbm, dsfont, xspace, cancel, mathtools). The delivered numbering is the unit's own. Assets: no retained span contains a figure environment, a graphics-inclusion command, a PDF-page inclusion, a table or a TikZ picture; no image file is copied into this package, the package declares no asset dependency and no image file is redistributed with it. Licence: version arXiv:2601.12165v1 of this article is distributed under the Creative Commons Attribution 4.0 International licence, as recorded in the arXiv licence metadata for this exact version and shown on the abstract page https://arxiv.org/abs/2601.12165v1. A full rights notice is delivered at licenses/arxiv-2601.12165-CC-BY-4.0.txt. Characters: every retained excerpt is pure ASCII, so the delivered engine, XeLaTeX with its default Latin Modern text font, reports no missing character.
\begin{lemma} \label{lm:partitions PsiAB}
	Let $(N,b) $ be a root and $\lambda, \mu\in \mathcal{M}_{b,N}^{(A,B)}$ such that 
	\begin{align*}
		\langle \Phi_\lambda, AB \Phi_\mu \rangle \neq 0.
	\end{align*}
	Then the renewal points of $\mu, \lambda$ strictly between the elementary blocks  $\max_q(A) $ and $ min_q(B) $ coincide.
	Furthermore, either:
	\begin{enumerate}
	\item $ \max_q(A)<  L_1 < L_2 <  \min_q(B) $ and every partition in $\mathcal{M}_{b,N}^{(A,B)}$ has a renewal point between the elementary blocks $\max_q(A)$ (excluding) and $ L_1$ (including) as well as  between $L_2$ (including) and $ \min_q(B)$ (excluding), or, if not, then
	\item there exists an elementary block longer than or equal to $ d_{AB} /3 $ with non-empty intersection with the orbitals in $ [ \max \supp A + d_{AB}/3 , \min \supp B - d_{AB}/3]$ which is separated from the rest by renewal points that are common to all partitions in $\mathcal{M}_{b,N}^{(A,B)}$.
	\end{enumerate}
\end{lemma}

\begin{proof}

	Let us first consider the case that the elementary blocks addressed by $ \max_q(A) $ and $ \min_q(B) $ are shorter than $ d_{AB}/3 $, i.e.
	$$
		b_{\max_q(A)}-b_{\max_q(A)-1}+q< d_{AB}/3, \qquad b_{\min_q(B)}-b_{\min_q(B)-1}+q< d_{AB}/3.
	$$
	This implies $\max_q(A)<L_1$ and $L_2<\min_q(B)$. In this situation, first suppose that
	 $\max_q(A)<L_1<L_2<\min_q(B)$. By construction, every element of $\mathcal{M}_{b,N}^{(A,B)}$ must then have one renewal point after the elementary block $\max_q(A) $ and before $L_1$ and one between $L_2 $ and $ \min_q(B)$. 
	 We define the number of particles, respectively the momentum generated by $ B = c_{L'}^* c_L $ by
	\begin{align*}
		\Delta_B \coloneqq \vert L'\vert - \vert L \vert, \qquad p_B \coloneqq \sum_{k\in L'} k - \sum_{k\in L} k, 
	\end{align*}
	and assume without loss of generality  that $\Delta_B\geq 0$, since otherwise we swap $\mu$ and $\lambda$ (effectively swapping $L$ and $L'$). 
	For any  pair of partitions $\lambda, \mu \in \mathcal{M}_{b,N}^{(A,B)}$, which we may identify with multisets through their occupation numbers, and which satisfy
	$
		\langle \Phi_\mu, c_{K'}^* c_K c_{L'}^* c_L \Phi_\lambda \rangle \neq 0 $, 
	we have 
	\begin{equation} \label{mu from lambda}
		\mu = (\lambda\setminus (L\cup K)) \cup (K'\cup L') . 
	\end{equation} 
	In particular, any entry $\mu_j $ on the elementary blocks from $ \max \mathrm{supp}(A)+1 $ to $ \min \mathrm{supp}(B)-1$ must correspond to some $\lambda_k$ (as those momenta are left unchanged by the action of $A$ and $B$). Let $s\leq L_1$ be a renewal point of $\mu$, i.e.
	\begin{equation} \label{renewal condition mu}
		\sum_{j=s+1}^N (\mu_j - (\lambda_b^{(q)})_j) = 0.
	\end{equation}   
	The relation \eqref{mu from lambda} implies
	\begin{equation} \label{momenta difference}
		\sum_{j=s+1}^N \mu_j - \sum_{j=s-\Delta_B+1}^N \lambda_j = -p_B.
	\end{equation}
	Indeed, \eqref{mu from lambda} tells us that we obtain $\mu$ from $\lambda$ by replacing the momenta in $L\cup K$ by the ones in $K'\cup L'$. Combining \eqref{renewal condition mu} and \eqref{momenta difference}, we obtain
	\begin{equation} \label{renewal equation mu}
		0 = \sum_{j=s+1}^N (\mu_j-(\lambda_b^{(q)})_j)
		= \sum_{j=s-\Delta_B+1}^N (\lambda_j-(\lambda_b^{(q)})_j) -p_B + \sum_{j=s-\Delta_B+1}^s (\lambda_b^{(q)})_j.
	\end{equation}
	As $\lambda \preceq \lambda_b^{(q)}(N)$, we hence conclude
	\begin{equation} \label{upper bound pB}
		p_B - \sum_{j=s-\Delta_B+1}^s (\lambda_b^{(q)})_j = \sum_{j=s-\Delta_B+1}^N (\lambda_j-(\lambda_b^{(q)})_j) \leq 0.
	\end{equation}
	Next, we pick a renewal point $\tilde{s}\geq L_2$ of $\lambda$. Using again \eqref{mu from lambda}, we obtain
	\begin{align*}
		\sum_{j=\tilde{s}+1}^N \lambda_j - \sum_{j=\tilde{s}+\Delta_B+1}^N \mu_j = p_B.
	\end{align*}
	As $\tilde{s}$ is a renewal point of $\lambda$, we get
	\begin{align*}
		0 = \sum_{j=\tilde{s}+1}^N (\lambda_j-(\lambda_b^{(q)})_j)
		= \sum_{j=\tilde{s}+\Delta_B+1} (\mu_j -(\lambda_b^{(q)})_j) + p_B -\sum_{j=\tilde{s}+1}^{\tilde{s}+\Delta_B} (\lambda_b^{(q)})_j
	\end{align*}
	and therefore, using $\mu \preceq \lambda_b^{(q)}(N)$, 
	\begin{equation} \label{lower bound pB}
		p-\sum_{j=\tilde{s}+1}^{\tilde{s}+\Delta_B} (\lambda_b^{(q)})_j = \sum_{j=\tilde{s}+\Delta_B+1}^N ((\lambda_b^{(q)})_j-\mu_j) \geq 0.
	\end{equation}
	Combining \eqref{upper bound pB} and \eqref{lower bound pB}, we arrive at
	\begin{equation} \label{bounds pB}
		\sum_{j=\tilde{s}+1}^{\tilde{s}+\Delta_B} (\lambda_b^{(q)})_j \leq p_B \leq \sum_{j=s-\Delta_B+1}^s (\lambda_b^{(q)})_j.
	\end{equation}
	However, $\tilde{s}\geq L_2>L_1\geq s$ and the entries of $\lambda_b^{(q)}(N)$ are strictly increasing. Thus, \eqref{bounds pB} is only possible if $p_B=0=\Delta_B$. However, in this case \eqref{momenta difference} implies that $s$ is also a renewal point of $\lambda$. Thus, $\lambda$ and $\mu$ share a common renewal point  on the elementry blocks $ \max_q(A)+1, \dots , \min_q(B)-1$. However, because the entries on those blocks must coincide, one obtains that all the renewal points on those blocks must coincide.
	
	If $L_1=L_2$, then $
		b_{L_1}-b_{L_1-1} +q  \geq d_{AB}/3 $, and hence the elementary block $ L_1 $ is separated from the rest by renewal points for all
	$\lambda\in \mathcal{M}_{b,N}^{(A,B)}$. Again, $\mu$ and $\lambda$ must share all their renewal points on the blocks between $ \max_q(A)$ and $  \min_q(B)$. 
	
	We now address the case that the elementary block $ \max_q(A) $ has length
	$$ b_{\max_q(A)}-b_{\max_q(A)-1}+q\geq d_{AB}/3. $$
	By definition of $\Psi_{b,N}^{(AB)}$, the elementary block $\max_q(A)$ is then separated by the rest through renewal points for all elements in $\mathcal{M}_{b,N}^{(A,B)}$. Since the entries of $\mu, \lambda$ in the on the orbitals $ \max \mathrm{supp}(A)+1, \dots ,  \min \mathrm{supp}(B)-1$ must coincide, this implies that all the renewal points on the blocks $\max_q(A), \dots , \min_q(B)-1$ must coincide.
	
	The case $ b_{\min_q(B)}-b_{\min_q(B)-1}+q\geq d_{AB}/3 $ works the same way.
\end{proof}

\end{document}
